\section{Algebraic and analytic conventions}\label{sec:preliminaries}

Throughout, $V$ has the hypotheses of Theorem~\ref{thm:main}. We collect the algebraic inputs that do not require positivity of all modules, and fix the operator conventions used in the estimates.

\subsection{Modules, duality and fusion}

By the regularity theorem of Abe--Buhl--Dong \cite[Theorem~4.5]{ABD}, every weak $V$-module is an algebraic direct sum of ordinary simple modules. There are finitely many simple isomorphism classes. Their lowest conformal weights are rational \cite[Theorem~11.3]{DLMModularInvariance}. These weights are not assumed nonnegative before a positive invariant form is known. In the grading-restricted category, each simple occurs with finite multiplicity: infinitely many copies would give an infinite-dimensional space at its lowest weight. Thus the categories used here have finite length.

The Huang--Lepowsky tensor theory supplies the fusion product $M\boxtimes N$ and its universal intertwining operator. For a target $K$, postcomposition identifies
\[
 \Hom_V(M\boxtimes N,K)
 \quad\text{with intertwining operators of type}\quad
 \binom{K}{M\ N}.
\]
Products and iterates of these operators have absolutely convergent matrix coefficients in their respective regions, and the fusing isomorphism identifies them. We use the positive-real convention
\[
 x_1>x_2>x_1-x_2>0
\]
for this identification, with real logarithms. The convergence and associativity statements include sums over intermediate conformal-energy projections \cite[Theorems~3.5 and~3.7, Remark~3.8]{HuangTensor}. Huang's rigidity and modularity theorems apply under precisely the present strong rationality hypotheses \cite[Theorems~1.1, 3.8 and~4.6]{HuangRigidity}. In particular the graded contragredient $M^\vee$ is a categorical dual. No positive form on $M$ is needed for these algebraic assertions.

We use ordinary, nonlogarithmic intertwining operators; the modules under consideration are semisimple with diagonalizable $L_0$. Their skew-symmetry and contragredient transformations are those of Frenkel--Huang--Lepowsky \cite{FHL}. The conjugate module $\bar M$ is the conjugate vector space with the action conjugated using $\Theta$. For an arbitrary module its \emph{Hermitian dual} is
\[
 M^{\mathrm h}=(\bar M)^\vee.
\]
The natural pairing is
\begin{equation}\label{eq:hermitian-pairing}
 [x,\lambda]_M=\overline{\lambda(\bar x)},
 \qquad x\in M,\quad\lambda\in M^{\mathrm h}.
\end{equation}
It is linear in $x$ and conjugate-linear in $\lambda$. This pairing has no positivity assertion. An invariant sesquilinear form on $M$ corresponds to a module map $M\to M^{\mathrm h}$.

All Hermitian inner products in the paper are linear in their first variable. A unitary module is a finite orthogonal sum of unitary simples: the grade-wise orthogonal complement of a submodule is invariant by the mode adjoint identities. A homomorphism between finite-length unitary modules extends to a bounded operator on their completions, and its adjoint is again a module homomorphism. To see the boundedness, decompose into simple types. On an irreducible type invariant positive forms differ by a positive scalar, by Schur's lemma; the remaining maps act on finite-dimensional multiplicity spaces.

The spectrum of $L_0$ on a unitary module is nonnegative. At a lowest vector of weight $h$, the relation $[L_1,L_{-1}]=2L_0$ gives $\|L_{-1}x\|^2=2h\|x\|^2$. Every weight in that simple module is $h+n$ for an integer $n\geq0$; finite orthogonal sums give the general case. Constants in our estimates may depend on the fixed module; uniformity over all modules is unnecessary.

\subsection{Modes, smooth vectors and closed operators}

For a homogeneous state $a$ of weight $d$, the \emph{physical modes} are
\begin{equation}\label{eq:physical-modes}
 a[m]=a_{(m+d-1)},\qquad
 [L_0,a[m]]=-m\,a[m].
\end{equation}
We use the same convention for intertwining fields, whose physical indices may lie in nonintegral translates of $\mathbb Z$. On an algebraic module, let $\Pi_\lambda$ denote the projection onto the $L_0$-eigenspace of energy $\lambda$ along the other eigenspaces. This definition does not require a Hermitian form. On a unitary completion, the same symbol denotes its orthogonal extension, and $P_N$ denotes the spectral projection onto energies at most $N$. We write
\[
 \mathcal H_M^\infty
 =\bigcap_{k\geq0}\Dom(1+L_0)^k
\]
for the smooth vectors, with their energy Sobolev seminorms. Their finite-energy subspace is the algebraic module $M$.

A state is \emph{quasi-primary} if $L_1a=0$, and \emph{primary} if $L_na=0$ for every $n>0$. A quasi-primary state of integral weight $d$ is called Hermitian when $(-1)^d\Theta a=a$. The invariant form then gives the level-wise adjoint identity
\begin{equation}\label{eq:hermitian-mode-adjoint}
 a[m]^*=a[-m].
\end{equation}
For a general state, the adjoint formula has the additional, finitely many $L_1$-descendant terms. We use \eqref{eq:hermitian-mode-adjoint} only with its stated hypothesis.

Polynomial mode bounds imply that smooth smearings and their formal adjoints map smooth vectors continuously to smooth vectors. Indeed Fourier coefficients of a smooth test decay faster than every power, and commuting $(1+L_0)^k$ across a mode costs only a polynomial in its energy shift. Finite-energy truncations converge in all these seminorms. Consequently the closures of a polynomially bounded smear from finite-energy vectors and from smooth vectors agree. We will repeatedly use this explicit core.

For a charged field, smearing is performed on an interval with a chosen angular lift; its physical normalization is $z^d\mathcal Y(a,z)$. Fractional powers use the branch specified by that lift. Changing the normalization by a smooth nowhere-zero phase merely changes the test function. At a reflected coordinate the conjugate branch is always used.

Positive-energy unitary Virasoro representations integrate to the centrally extended covering group of circle diffeomorphisms and satisfy the usual linear energy bounds \cite[Proposition~3.2, Remark~3.3 and Theorem~3.4]{CKLW}; see also \cite[Theorem~5.2.1 and Proposition~5.3.1]{ToledanoLaredo}. The resulting M\"obius and diffeomorphism operators preserve the smooth space. We use their fixed conformal normalization; in particular the Virasoro relation $[L_1,L_{-1}]=2L_0$ fixes the additive constant in $L_0$. Annulus actions and their Sobolev properties will be introduced only after the local net has been constructed.

For a closed densely defined operator $T$, its graph is a closed subspace of the direct sum of its source and target Hilbert spaces. The entries of the orthogonal graph projection are bounded operators. Graph projections allow us to state precisely that a normal representation carries one closed field to another, or that a bounded projection reduces a closed field. Agreement on a dense set alone will not be used as a substitute for agreement of closed graphs.

\subsection{Creation and the adjoint convention}

For a unitary module $M$, the creation intertwining operator is
\begin{equation}\label{eq:creation}
 \mathcal Y_M^+(a,z)v=e^{zL_{-1}}Y^M(v,-z)a,
 \qquad a\in M,\quad v\in V.
\end{equation}
Its type is $\binom{M}{M\ V}$. We use Tener's adjoint convention \cite[Section~2 and Proposition~3.10]{TenerFusion}; its adjoint, the annihilation operator, is
\[
 \mathcal Y_M^-=(\mathcal Y_M^+)^\dagger,
 \qquad\text{of type }\binom{V}{\bar M\ M}.
\]
For positive $r$, put
\begin{equation}\label{eq:theta-convention}
 \theta_r a=e^{rL_1}e^{i\pi L_0}r^{-2L_0}\bar a\in\bar M.
\end{equation}
For complex coordinates the logarithms follow the chosen angular lift. The reflection formula uses $\theta_{\bar w}$ and the coordinate $\bar w^{-1}$. These conventions fix the phase of every transport form below, including when the lowest weight of $M$ is nonintegral. The comparison with Gui's canonical fusion normalization is made explicitly in the positivity argument.

Tener's Propositions~3.10--3.11 \cite{TenerFusion} give two inputs we will need before general fusion positivity. First, creation satisfies the vacuum-source transport identity. Second, for unitary charge, source and target, an intertwining operator applied to finite-energy vectors has Hilbert-space values on $0<|z|<1$, holomorphic on each branch. The latter conclusion uses regularity and convergence of products with the adjoint intertwiner. It does not require an energy bound or a net representation on the source. We will distinguish this Hilbert-valued convergence from the operator bounds that we prove later.

\subsection{The scalar case}

If $V=\mathbb C\Omega$, the conformal vector is zero, every grading-restricted module is finite dimensional in degree zero, and all fields are constant scalar fields. The net on the one-dimensional vacuum space is the scalar net. It is irreducible, split and strongly additive, and its two-interval index is one. Its additive dualizable representation category is the category of finite-dimensional Hilbert spaces. Creation, fusion and CWX integration are the usual tensor product over $\mathbb C$, with the flip braiding. Thus Theorem~\ref{thm:main} holds in this case.

We assume $V\neq\mathbb C\Omega$ henceforth. Its central charge is positive: $\|\nu\|^2=c/2$, so $c=0$ would force $\nu=0$, and then CFT type would force $V=V_0=\mathbb C\Omega$. The nontriviality assumption is used in the flip-orbifold argument and in the usual conformal-net sector theory.
